\chapter{Chapter08 — Composed Majority Influence}
%%% generated by the blueprint pipeline from FABL.Chapter08.ComposedMajorityInfluence
\section{Overview}

Declarations of the Lean module \texttt{FABL.Chapter08.ComposedMajorityInfluence}.

\section{Declarations}

\begin{theorem}[Influence of composed majority]
\lean{FABL.totalInfluence_composedMajority}
\label{FABL.totalInfluence_composedMajority}
\leanok
\uses{FABL.BooleanFunction.toReal, FABL.Sign, FABL.biasSigma, FABL.biasSigma_sq, FABL.biasVarianceScale, FABL.biasedMinusProbabilityPolynomial, FABL.biasedSignPMF, FABL.biasedSignPMF_one_half_eq_uniform, FABL.composedMajority, FABL.composedMajority_monotone, FABL.composedMajority_probability_formula, FABL.hasDerivAt_biasedMinusProbabilityPolynomial_eq_totalInfluence, FABL.majority, FABL.majority_monotone, FABL.productTotalInfluence_uniformSign_eq_totalInfluence, FABL.totalInfluence, FABL.uniformPMF}
\difficulty{6}
For every natural number $m$, the total influence of the composed majority function
equals $1$ plus one half of the total influence of majority on sign vectors of dimension
$2m+1$.
\end{theorem}

\begin{definition}[Composed majority influence error]
\lean{FABL.composedMajorityInfluenceError}
\label{FABL.composedMajorityInfluenceError}
\leanok
\uses{FABL.BooleanFunction.toReal, FABL.composedMajority, FABL.totalInfluence}
For each $m\in\mathbb N$, define the composed majority influence error as the total
influence of the composed majority function on $m$ coordinates, minus
$1+\frac{1}{2}\sqrt{\frac{2}{\pi}}\sqrt{2m+1}$.
\end{definition}

\begin{theorem}[Composed majority influence error identity]
\lean{FABL.composedMajorityInfluenceError_eq}
\label{FABL.composedMajorityInfluenceError_eq}
\leanok
\uses{FABL.composedMajorityInfluenceError, FABL.oddMajorityTotalInfluenceError, FABL.totalInfluence_composedMajority, FABL.totalInfluence_majority_odd_eq_main_add_error}
\difficulty{2}
For every integer $m\ge 0$, the total influence of the composed majority on the sign
cube of dimension $2+(2m+1)$, minus $1+\frac12\sqrt{2/\pi}\sqrt{2m+1}$, equals one half
of the total influence of majority on the sign cube of dimension $2m+1$, minus
$\sqrt{2/\pi}\sqrt{2m+1}$.
\end{theorem}

\begin{theorem}[Bounds for composed majority influence error]
\lean{FABL.composedMajorityInfluenceError_mem_Icc}
\label{FABL.composedMajorityInfluenceError_mem_Icc}
\leanok
\uses{FABL.composedMajorityInfluenceError, FABL.composedMajorityInfluenceError_eq, FABL.oddMajorityTotalInfluenceError, FABL.oddMajorityTotalInfluenceError_mem_Icc}
\difficulty{4}
For every integer $m\ge 0$, the total influence of the composed majority on the sign
cube of dimension $2+(2m+1)$, minus $1+\frac12\sqrt{2/\pi}\sqrt{2m+1}$, lies in the
interval $[0,\,1/(2\sqrt{2m+1})]$.
\end{theorem}

\begin{theorem}[Composed majority influence identity]
\lean{FABL.totalInfluence_composedMajority_eq_main_add_error}
\label{FABL.totalInfluence_composedMajority_eq_main_add_error}
\leanok
\uses{FABL.BooleanFunction.toReal, FABL.composedMajority, FABL.composedMajorityInfluenceError, FABL.totalInfluence}
\difficulty{1}
For every natural number $m$, the total influence of the composed majority function on
sign vectors of dimension $2+(2m+1)$ equals
$1+\frac{1}{2}\sqrt{\frac{2}{\pi}}\sqrt{2m+1}$ plus its influence error, where the
influence error is the total influence minus
$1+\frac{1}{2}\sqrt{\frac{2}{\pi}}\sqrt{2m+1}$.
\end{theorem}

\begin{theorem}[Lower bound for the Gaussian constant]
\lean{FABL.one_half_le_sqrt_two_div_pi}
\label{FABL.one_half_le_sqrt_two_div_pi}
\leanok
\difficulty{3}
The real number $\sqrt{2/\pi}$ is at least $1/2$.
\end{theorem}

\begin{theorem}[A bound involving pi]
\lean{FABL.sqrt_two_div_pi_le_one}
\label{FABL.sqrt_two_div_pi_le_one}
\leanok
\difficulty{3}
The square root of $2/\pi$ is at most $1$.
\end{theorem}

\begin{theorem}[Composed Majority Dimension Bound]
\lean{FABL.sqrt_composedMajority_dimension_le}
\label{FABL.sqrt_composedMajority_dimension_le}
\leanok
\difficulty{4}
For every natural number $m$, $\sqrt{2m+3}\le 2\sqrt{2m+1}$.
\end{theorem}

\begin{theorem}[Majority Tail Dimension Bound]
\lean{FABL.sqrt_majorityTail_le_composedMajority_dimension}
\label{FABL.sqrt_majorityTail_le_composedMajority_dimension}
\leanok
\difficulty{3}
For every natural number $m$, $\sqrt{2m+1}\le \sqrt{2m+3}$.
\end{theorem}

\begin{theorem}[Lower bound for composed majority influence]
\lean{FABL.totalInfluence_composedMajority_lower}
\label{FABL.totalInfluence_composedMajority_lower}
\leanok
\uses{FABL.BooleanFunction.toReal, FABL.composedMajority, FABL.composedMajorityInfluenceError, FABL.composedMajorityInfluenceError_mem_Icc, FABL.one_half_le_sqrt_two_div_pi, FABL.sqrt_composedMajority_dimension_le, FABL.totalInfluence, FABL.totalInfluence_composedMajority_eq_main_add_error}
\difficulty{4}
For every natural number $m$, the total influence of the composed majority function on
sign vectors of dimension $2+(2m+1)$ is at least $\frac{1}{8}\sqrt{2m+3}$.
\end{theorem}

\begin{theorem}[Upper bound for composed majority influence]
\lean{FABL.totalInfluence_composedMajority_upper}
\label{FABL.totalInfluence_composedMajority_upper}
\leanok
\uses{FABL.BooleanFunction.toReal, FABL.composedMajority, FABL.composedMajorityInfluenceError, FABL.composedMajorityInfluenceError_mem_Icc, FABL.sqrt_majorityTail_le_composedMajority_dimension, FABL.sqrt_two_div_pi_le_one, FABL.totalInfluence, FABL.totalInfluence_composedMajority_eq_main_add_error}
\difficulty{5}
For every natural number $m$, the total influence of the composed majority function on
sign vectors of dimension $2+(2m+1)$ is at most $2\sqrt{2m+3}$.
\end{theorem}

\begin{theorem}[Asymptotic influence growth]
\lean{FABL.totalInfluence_composedMajority_isTheta_sqrt_dimension}
\label{FABL.totalInfluence_composedMajority_isTheta_sqrt_dimension}
\leanok
\uses{FABL.BooleanFunction.toReal, FABL.composedMajority, FABL.totalInfluence, FABL.totalInfluence_composedMajority_lower, FABL.totalInfluence_composedMajority_upper, FABL.totalInfluence_nonneg}
\difficulty{4}
As $m$ tends to infinity through the natural numbers, the total influence of the
composed majority function is asymptotic in order to $\sqrt{2m+3}$; that is, it is
bounded above and below by positive constant multiples of this quantity for all
sufficiently large $m$.
\end{theorem}

\begin{definition}[Composed majority arity]
\lean{FABL.composedMajorityArity}
\statementsource{Analysis-of-Boolean-Functions-by-Ryan-ODonnell}{pdf-ffbe9c7b3306-p044-b007}
\label{FABL.composedMajorityArity}
\leanok
For every natural number $m$, the composed-majority arity is $2m+3$.
\end{definition}

\begin{definition}[Composed majority family]
\lean{FABL.composedMajorityFamily}
\statementsource{Analysis-of-Boolean-Functions-by-Ryan-ODonnell}{
  pdf-ffbe9c7b3306-p044-b007,
  pdf-ffbe9c7b3306-p239-b004
}
\label{FABL.composedMajorityFamily}
\leanok
\uses{FABL.BooleanFunction, FABL.composedMajority, FABL.composedMajorityArity}
For each natural number $m$, let $F_m$ be the composed-majority Boolean function on the
sign cube of dimension $2m+3$, taking values in $\{ -1,1 \}$.
\end{definition}

\begin{theorem}[Monotonicity of composed majority]
\lean{FABL.composedMajorityFamily_monotone}
\label{FABL.composedMajorityFamily_monotone}
\leanok
\uses{FABL.composedMajorityFamily, FABL.composedMajority_monotone}
\difficulty{1}
For every integer $m\ge 0$, the composed majority function on the sign cube of dimension
$2+(2m+1)$ is monotone: if $x_i\le y_i$ in every coordinate, then its value at $x$ is at
most its value at $y$, where $-1<1$.
\end{theorem}

\begin{theorem}[Nonconstancy of composed majority]
\lean{FABL.composedMajorityFamily_isNonconstant}
\label{FABL.composedMajorityFamily_isNonconstant}
\leanok
\uses{FABL.composedMajorityFamily, FABL.composedMajority_isNonconstant}
\difficulty{1}
For every natural number $m$, the composed majority function on $2+(2m+1)$ sign
coordinates is nonconstant.
\end{theorem}

\begin{theorem}[Normalized width of composed majority]
\lean{FABL.normalizedThresholdWidth_composedMajority_eq_thresholdWidth}
\label{FABL.normalizedThresholdWidth_composedMajority_eq_thresholdWidth}
\leanok
\uses{FABL.biasVarianceScale, FABL.composedMajority, FABL.composedMajorityArity, FABL.composedMajorityFamily, FABL.composedMajorityFamily_isNonconstant, FABL.composedMajorityFamily_monotone, FABL.composedMajority_isNonconstant, FABL.composedMajority_monotone, FABL.criticalProbability_composedMajority, FABL.normalizedThresholdWidth, FABL.thresholdWidth}
\difficulty{3}
For every $\epsilon\in(0,1/2)$ and every natural number $m$, the normalized threshold
width of the composed-majority family at index $m$ equals the threshold width of
composed majority on $2+(2m+1)$ inputs, with both widths taken at $\epsilon$.
\end{theorem}

\begin{definition}[Composed majority threshold probe]
\lean{FABL.composedMajorityThresholdProbe}
\label{FABL.composedMajorityThresholdProbe}
\leanok
For each $\epsilon\in\mathbb R$, define the composed majority threshold probe to be
$\frac{\sqrt{\epsilon}+1/2}{2}$.
\end{definition}

\begin{theorem}[Threshold probe lies in an interval]
\lean{FABL.composedMajorityThresholdProbe_mem_Ioo}
\label{FABL.composedMajorityThresholdProbe_mem_Ioo}
\leanok
\uses{FABL.composedMajorityThresholdProbe}
\difficulty{4}
If $0<\epsilon<1/4$, then $(\sqrt{\epsilon}+1/2)/2$ lies strictly between $0$ and $1/2$.
\end{theorem}

\begin{theorem}[Complement of the majority threshold probe]
\lean{FABL.one_sub_composedMajorityThresholdProbe_mem_Ioo}
\label{FABL.one_sub_composedMajorityThresholdProbe_mem_Ioo}
\leanok
\uses{FABL.composedMajorityThresholdProbe, FABL.composedMajorityThresholdProbe_mem_Ioo}
\difficulty{4}
Let $0<\epsilon<1/4$, and define the composed majority threshold probe by
$T(\epsilon)=(\sqrt{\epsilon}+1/2)/2$. Then $1-T(\epsilon)$ lies strictly between $1/2$
and $1$.
\end{theorem}

\begin{theorem}[Threshold probe squared exceeds epsilon]
\lean{FABL.epsilon_lt_composedMajorityThresholdProbe_sq}
\label{FABL.epsilon_lt_composedMajorityThresholdProbe_sq}
\leanok
\uses{FABL.composedMajorityThresholdProbe}
\difficulty{4}
For every real number $\epsilon$ with $0<\epsilon<\tfrac14$, one has
$\epsilon<\left(\frac{\sqrt{\epsilon}+\tfrac12}{2}\right)^2$.
\end{theorem}

\begin{theorem}[A composed threshold identity]
\lean{FABL.one_sub_two_mul_composedMajorityThresholdProbe}
\label{FABL.one_sub_two_mul_composedMajorityThresholdProbe}
\leanok
\uses{FABL.composedMajorityThresholdProbe}
\difficulty{3}
For every real number $\epsilon$, if $t=(\sqrt{\epsilon}+\frac12)/2$, then
$1-2t=\frac12-\sqrt{\epsilon}$.
\end{theorem}

\begin{theorem}[Interval inclusion for epsilon]
\lean{FABL.epsilon_mem_Ioo_half_of_mem_Ioo_quarter}
\label{FABL.epsilon_mem_Ioo_half_of_mem_Ioo_quarter}
\leanok
\difficulty{1}
Every real number $\epsilon$ with $0<\epsilon<\tfrac14$ also satisfies
$0<\epsilon<\tfrac12$.
\end{theorem}

\begin{theorem}[Composed majority normalized width lower bound]
\lean{FABL.eventually_composedMajority_normalizedThresholdWidth_gt}
\label{FABL.eventually_composedMajority_normalizedThresholdWidth_gt}
\leanok
\uses{FABL.biasedMinusProbabilityPolynomial, FABL.biasedMinusProbabilityPolynomial_thresholdLowerParameter, FABL.biasedMinusProbabilityPolynomial_thresholdUpperParameter, FABL.biasedProbabilityLevel_mem_Ioo, FABL.composedMajority, FABL.composedMajorityArity, FABL.composedMajorityFamily, FABL.composedMajorityFamily_isNonconstant, FABL.composedMajorityFamily_monotone, FABL.composedMajorityThresholdProbe, FABL.composedMajorityThresholdProbe_mem_Ioo, FABL.composedMajority_isNonconstant, FABL.composedMajority_monotone, FABL.epsilon_lt_composedMajorityThresholdProbe_sq, FABL.epsilon_mem_Ioo_half_of_mem_Ioo_quarter, FABL.normalizedThresholdWidth, FABL.normalizedThresholdWidth_composedMajority_eq_thresholdWidth, FABL.one_sub_composedMajorityThresholdProbe_mem_Ioo, FABL.one_sub_two_mul_composedMajorityThresholdProbe, FABL.strictMonoOn_biasedMinusProbabilityPolynomial, FABL.tendsto_composedMajority_probability_of_half_lt, FABL.tendsto_composedMajority_probability_of_lt_half, FABL.thresholdLowerParameter, FABL.thresholdUpperParameter}
\difficulty{6}
For every $\epsilon\in(0,1/4)$, for all sufficiently large natural numbers $m$, the
normalized threshold width of the composed majority family at parameter $\epsilon$ is
greater than $1/2-\sqrt{\epsilon}$.
\end{theorem}

\begin{theorem}[Composed majority normalized threshold width not tendsto zero]
\lean{FABL.composedMajority_normalizedThresholdWidth_not_tendsto_zero}
\statementsource{Analysis-of-Boolean-Functions-by-Ryan-ODonnell}{
  pdf-ffbe9c7b3306-p238-b009,
  pdf-ffbe9c7b3306-p238-b010,
  pdf-ffbe9c7b3306-p239-b007
}
\label{FABL.composedMajority_normalizedThresholdWidth_not_tendsto_zero}
\leanok
\uses{FABL.composedMajorityArity, FABL.composedMajorityFamily, FABL.composedMajorityFamily_isNonconstant, FABL.composedMajorityFamily_monotone, FABL.epsilon_mem_Ioo_half_of_mem_Ioo_quarter, FABL.eventually_composedMajority_normalizedThresholdWidth_gt, FABL.normalizedThresholdWidth}
\difficulty{4}
% dataset-pipeline: stub -- prose pending build_informal_statements.py
Statement pending regeneration.
\end{theorem}

\begin{theorem}[Composed majority has a coarse threshold]
\lean{FABL.composedMajority_hasCoarseThreshold}
\label{FABL.composedMajority_hasCoarseThreshold}
\leanok
\uses{FABL.HasCoarseThreshold, FABL.composedMajorityArity, FABL.composedMajorityFamily, FABL.composedMajorityFamily_isNonconstant, FABL.composedMajorityFamily_monotone, FABL.composedMajority_normalizedThresholdWidth_not_tendsto_zero, FABL.epsilon_mem_Ioo_half_of_mem_Ioo_quarter}
\difficulty{1}
For every $0<\epsilon<1/4$, consider the sequence of composed majority functions whose
dimension at index $k$ is $2+(2k+1)$. For each function, let its lower and upper
threshold parameters be the unique biases in $(0,1)$ at which its probability of output
$-1$ under independent biased signs is respectively $\epsilon$ and $1-\epsilon$; let its
critical bias be the unique bias at which that probability is $1/2$. The sequence
obtained by dividing the difference of the upper and lower threshold parameters by
$4p(1-p)$ at the critical bias $p$ does not converge to $0$ as $k\to\infty$.
\end{theorem}
